\section*{Discussion}
This report focuses on not a decisive evaluation of the generalized applicability of this algorithm-collection, but the case-specific use on Xenium and Spleen data. The evaluation is limited in its merit, chiefly rooted in the necessity of a ground-truth and its apparent quality. The custom-pipeline, inevitably developed in tandem, though fully functional, has yet to be tested by unfamiliar users.\\
The reference-based analysis, under consideration of the visual interpretation, showed a clear dominance of \textsc{XOA}. The prerequisite to call upon the human ... examplifies the (trust) we can (put into) reference-based evaluations. 
- obvious limitations of the algortihms AND the evaluation methods.\\
to consider as future development: GANs (generative adversarial networks) to generate known "spleen" images. different approach: using statistical modeling to generate three-dimensional "spleen" tissue with known location and type of cell --- this could be supported by 2p-microscopy images? \\
XOA was trained on xenium-based manual segmentation --- thus my own segmentation is probably the closest to what others drew\\
the f1 DT value/calculation has to be scrutinised --- it appears untrustworthy\\
the PQ value has very very low values. why? \\
as seen, no solitary value should ever be used to quantify the quality of a prediction. instead a collection of different values must be considered.\\
the approach of CSE could be adapted to include the information of SRT for the protein/cell type analysis, since it obviously contains more information than the protein stains of Xenium.\\
cross eval is based on the expectation, that there are multiple alogrithms that perform well. also it should be performed in tandem with a typical reference based analysis to identify candidates that one would expect to perform well, that then can be verified to explain together a similar segmentation. this is a rudimentary approach as the probability one. the threshold should either be kept low at .5 or set to a value inferred from the gt-based analysis.\\
the transcript tables beign in pixel format might impede on prosegs performance..\\
the cross evalation can, without a gt, identify similarities in the predictions.\\ 
(According to Davide Chicco and Giuseppe Jurman, the F1 score is less truthful and informative than the Matthews correlation coefficient (MCC) in binary evaluation classification \href{https://doi.org/10.1186%2Fs12864-019-6413-7}{doi})\\
This was evident to the CSE developers, as they write in their paper. The analysis works best in tissues with heterogenous expression and cell types, with a multitude of markers. Additionally the cell-shape: it is not penalised when nonsensical shapes are proposed.\\